\chapter{Waring's Problem}
\begin{center}\itshape
Every number is a sum of four squares, of nine cubes,\\
of nineteen fourth powers, ``and so on''. For large $s$\\
the circle method counts the representations.
\end{center}

Edward Waring, Lucasian professor at Cambridge, stated in his \emph{Meditationes Algebraicae} of 1770 that every natural number is a sum of four squares, of nine cubes, of nineteen fourth powers, ``and so on''. Lagrange proved the four squares the same year. The cubes waited until 1909, when Wieferich showed that nine suffice, and in the same year Hilbert proved that for every $k$ some finite number $g(k)$ of $k$-th powers suffices. Hilbert's proof gave no reasonable bound.

The preceding chapter wrote $p(n)$ as an integral over the circle. The same idea counts sums of powers. In 1920 Hardy and Littlewood, in the first of their papers \emph{Partitio Numerorum}, wrote the number of representations of $n$ as a sum of $s$ $k$-th powers as a Fourier coefficient of the $s$-th power of a Weyl sum. The circle splits into major arcs, where the Weyl sum is computed, and minor arcs, where Weyl's inequality makes it small. In 1938 Hua Loo-Keng, who had taught himself mathematics in his father's shop in Jintan and had come to Cambridge in 1936 to work with Hardy, found a mean value estimate that makes the asymptotics work for $2^k+1$ powers.

This chapter proves Weyl's inequality for degree $k$ and Hua's lemma, both by Weyl's differencing, and the asymptotic formula for every $s\ge 2^k+1$. The singular series is positive for $s\ge 4k$, and therefore $G(k)\le 2^k+1$. Volume~14 proves the companion theorem for three primes.

\begin{reminder}[from Volumes~6 and this book]
\emph{Geometric sums.} For an interval $I$, $\bigl|\sum_{n\in I}e(\beta n)\bigr|\le\min\bigl(|I|+1,\frac1{2\|\beta\|}\bigr)$.
\emph{Sums of minima} (Volume~6). If $\bigl|\alpha-\frac aq\bigr|\le\frac1{q^2}$ with $(a,q)=1$, then $\sum_{1\le x\le X}\min\bigl(N,\frac1{\|\alpha x\|}\bigr)\le\bigl(\frac Xq+1\bigr)(4N+2q+2q\log q)$. For every $\varepsilon>0$, $\tau(m)\ll_\varepsilon m^\varepsilon$.
\emph{The Beta integral} (Chapter~7). $\int_0^1 t^{x-1}(1-t)^{y-1}\dd t=\frac{\Gamma(x)\Gamma(y)}{\Gamma(x+y)}$.
\emph{Fourier inversion} (Chapter~10). If $h\in L^1(\R)$ is continuous and $\hat h\in L^1$, then $h(x)=\int\hat h(\xi)e(x\xi)\dd\xi$ for every $x$.
\end{reminder}

\begin{lemma}[Dirichlet]\label{lem:dirichlet}
For real $\alpha$ and $Q\ge1$ there are coprime integers $a$ and $q$ with $1\le q\le Q$ and $\bigl|\alpha-\frac aq\bigr|\le\frac1{qQ}$.
\end{lemma}
\begin{proof}
Let $n=\lfloor Q\rfloor$. Cut $[0,1]$ into $n+1$ closed intervals of length $\frac1{n+1}$. The $n+2$ numbers $\{j\alpha\}$, $0\le j\le n$, and $1$ lie in them, so two lie in the same one; their difference gives integers $1\le q\le n$ and $a$ with $|q\alpha-a|\le\frac1{n+1}<\frac1Q$. Dividing $a$ and $q$ by their gcd only decreases $q$, and keeps $\bigl|\alpha-\frac aq\bigr|\le\frac1{qQ}$.
\end{proof}

\section{Weyl's inequality and Hua's lemma}

For Waring's problem the generating function is the Weyl sum
\[
f(\alpha)=\sum_{1\le x\le P}e(\alpha x^k),\qquad k\ge2,\ P\ge1,
\]
and the tool is differencing. For a function $\psi$ on the integers and $h\in\Z$, let $\Delta_h\psi(x)=\psi(x+h)-\psi(x)$, and for $\mathbf h=(h_1,\dots,h_j)$ let $\Delta_{\mathbf h}=\Delta_{h_1}\cdots\Delta_{h_j}$. Expanding, $\Delta_{\mathbf h}\psi(x)=\sum_{S\subseteq\{1,\dots,j\}}(-1)^{j-|S|}\psi(x+h_S)$ with $h_S=\sum_{i\in S}h_i$. For an interval $I\subseteq[1,P]$ let $T(\psi,I)=\sum_{x\in I\cap\Z}e(\psi(x))$, let $H=\Z\cap\left]-P,P\right[$, and let $I_{\mathbf h}$ be the set of $x$ with $x+h_S\in I$ for all $S$, an interval contained in $I$.

\begin{lemma}[Weyl differencing]\label{lem:weyldiff}
For $j\ge1$,
\[
|T(\psi,I)|^{2^j}\le(2P)^{2^j-j-1}\sum_{\mathbf h\in H^j}T(\Delta_{\mathbf h}\psi,I_{\mathbf h}),
\]
and the right side is real and non-negative.
\end{lemma}
\begin{proof}
For any interval $J\subseteq[1,P]$ and function $\phi$, writing $x=y+h$,
\[
|T(\phi,J)|^2=\sum_{x,y\in J}e(\phi(x)-\phi(y))=\sum_{h\in H}T(\Delta_h\phi,J\cap(J-h)).\tag{$\ast$}
\]
For $j=1$ this is the claim. Suppose the claim holds for $j$. By $(\ast)$ with $\phi=\Delta_{\mathbf h}\psi$ and $J=I_{\mathbf h}$, summed over $\mathbf h\in H^j$, the right side for $j+1$ is $(2P)^{2^{j+1}-j-2}\sum_{\mathbf h\in H^j}|T(\Delta_{\mathbf h}\psi,I_{\mathbf h})|^2\ge0$; and the right side for $j$ is $\sum_{\mathbf h\in H^{j-1}}|\cdots|^2\ge0$ by the same identity (for $j=1$ it is $|T(\psi,I)|^2$). Squaring the inequality for $j$ and applying Cauchy--Schwarz over the at most $(2P)^j$ vectors $\mathbf h$,
\[
|T(\psi,I)|^{2^{j+1}}\le(2P)^{2^{j+1}-2j-2}(2P)^j\sum_{\mathbf h\in H^j}|T(\Delta_{\mathbf h}\psi,I_{\mathbf h})|^2,
\]
which is the claim for $j+1$.
\end{proof}

For $\psi(x)=\alpha x^k$ the differences have a simple shape. For $j\le k$ and $x\in I_{\mathbf h}$,
\[
\Delta_{\mathbf h}x^k=h_1\cdots h_j\,p_j(x;\mathbf h),\qquad p_j(x;\mathbf h)=\int_{[0,1]^j}\frac{k!}{(k-j)!}(x+\mathbf t\cdot\mathbf h)^{k-j}\dd\mathbf t,\tag{5.1}
\]
by the fundamental theorem of calculus applied in each variable $t_i$ in turn. The point $x+\mathbf t\cdot\mathbf h$ is multi-affine in $\mathbf t$, so it lies in the convex hull of the points $x+h_S$, all in $[1,P]$; hence $p_j(x;\mathbf h)\ge\frac{k!}{(k-j)!}\ge1$, and for $j<k$, $p_j$ is strictly increasing in $x$. For $j=k-1$, $p_{k-1}(x;\mathbf h)=k!\,(x+\frac{h_1+\cdots+h_{k-1}}2)$.

\begin{theorem}[Weyl's inequality]\label{thm:weylk}
Let $k\ge2$ and $\varepsilon>0$. If $|\alpha-\frac aq|\le\frac1{q^2}$ with $(a,q)=1$, then
\[
|f(\alpha)|\ll_{k,\varepsilon}P^{1+\varepsilon}\Bigl(\frac1q+\frac1P+\frac q{P^k}\Bigr)^{2^{1-k}}.
\]
\end{theorem}
\begin{proof}
We may assume $q\le P^k$, since otherwise the bound exceeds $P$. Apply Lemma~\ref{lem:weyldiff} with $j=k-1$ and $I=[1,P]$. By (5.1), $\Delta_{\mathbf h}(\alpha x^k)$ is linear in $x$ with slope $\alpha k!h_1\cdots h_{k-1}$, so the reminder bounds each inner sum by $\min(P+1,\frac1{2\|\alpha k!h_1\cdots h_{k-1}\|})$, and
\[
|f(\alpha)|^{2^{k-1}}\le(2P)^{2^{k-1}-k}\sum_{\mathbf h\in H^{k-1}}\min\Bigl(2P,\frac1{\|\alpha k!h_1\cdots h_{k-1}\|}\Bigr).
\]
The $\mathbf h$ with some $h_i=0$ contribute at most $(k-1)(2P)^{k-2}\cdot2P\ll P^{k-1}$. For the others, $m=k!|h_1\cdots h_{k-1}|$ lies in $[1,k!P^{k-1}]$, and a given $m$ comes from at most $2^{k-1}\tau(m)^{k-1}\ll P^\varepsilon$ vectors $\mathbf h$, since each $|h_i|$ divides $m$. With $X=k!P^{k-1}$, the second reminder bounds the sum over $m$ by
\[
\Bigl(\frac Xq+1\Bigr)(8P+2q+2q\log q)\ll\Bigl(\frac{P^k}q+P^{k-1}+P+q\Bigr)\log P\ll P^k\Bigl(\frac1q+\frac1P+\frac q{P^k}\Bigr)\log P.
\]
Altogether $|f(\alpha)|^{2^{k-1}}\ll P^{2^{k-1}-k}\cdot P^{k+2\varepsilon}(\frac1q+\frac1P+\frac q{P^k})$; take the $2^{k-1}$-th root and rename $\varepsilon$.
\end{proof}

For $k=2$ the saving is a square root. The exponent $2^{1-k}$ is the price of $k-1$ applications of Cauchy--Schwarz; Vinogradov's mean value theorem, recorded in Volume~14, improves it to about $\frac1{k^2}$, but we shall not need that.

\begin{theorem}[Hua, 1938]\label{thm:hua}
For $1\le j\le k$ and $\varepsilon>0$, $\displaystyle\int_0^1|f(\alpha)|^{2^j}\dd\alpha\ll_{k,\varepsilon}P^{2^j-j+\varepsilon}$.
\end{theorem}
\begin{proof}
For $j=1$, Parseval gives $\int_0^1|f|^2=\lfloor P\rfloor\le P$. Suppose the bound holds for some $j\le k-1$. Expanding,
\[
|f(\alpha)|^{2^j}=f(\alpha)^{2^{j-1}}\overline{f(\alpha)}^{2^{j-1}}=\sum_m\rho(m)e(\alpha m),
\]
where $\rho(m)\ge0$ counts the solutions of $y_1^k+\cdots+y_{2^{j-1}}^k-z_1^k-\cdots-z_{2^{j-1}}^k=m$ in integers of $[1,P]$; and $\rho(m)\le\rho(0)=\int_0^1|f|^{2^j}$ by Cauchy--Schwarz, $\sum_m\rho(m)=\lfloor P\rfloor^{2^j}$. Lemma~\ref{lem:weyldiff} and (5.1) give
\[
|f(\alpha)|^{2^j}\le(2P)^{2^j-j-1}\sum_{\mathbf h\in H^j}\sum_{x\in I_{\mathbf h}}e\bigl(\alpha D(x;\mathbf h)\bigr),\qquad D(x;\mathbf h)=h_1\cdots h_jp_j(x;\mathbf h),
\]
with a non-negative right side. Multiplying by $|f(\alpha)|^{2^j}$ and integrating, since $\int_0^1e(\alpha D)|f|^{2^j}=\rho(-D)$,
\[
\int_0^1|f|^{2^{j+1}}\le(2P)^{2^j-j-1}\sum_{\mathbf h\in H^j}\sum_{x\in I_{\mathbf h}}\rho(-D(x;\mathbf h)).
\]
If some $h_i=0$, then $D=0$; there are at most $j(2P)^{j-1}P$ such pairs $(\mathbf h,x)$, each contributing $\rho(0)\ll P^{2^j-j+\varepsilon}$, in total $O(P^{2^j+\varepsilon})$. If no $h_i$ is $0$, then $D\ne0$ and $|D|\le k!P^k$. A given value $m=D$ determines the $h_i$ up to $(2\tau(|m|))^j\ll P^\varepsilon$ choices, since each divides $m$; then $p_j=\frac m{h_1\cdots h_j}$, which determines $x$ because $p_j$ is strictly increasing in $x$. So these pairs contribute at most $O(P^\varepsilon)\sum_m\rho(-m)\ll P^{2^j+\varepsilon}$. Hence $\int_0^1|f|^{2^{j+1}}\ll P^{2^j-j-1}\cdot P^{2^j+\varepsilon}=P^{2^{j+1}-(j+1)+\varepsilon}$.
\end{proof}

For $j=k$ the lemma says that $\int|f|^{2^k}\ll P^{2^k-k+\varepsilon}$, which is the number of trivial solutions $P^{2^{k-1}}$ only when $k=2$, but is always within $P^\varepsilon$ of the lower bound $P^{2^k-k}$: the $P^{2^k}$ tuples of $2^k$ variables produce values in a range of length $O(P^k)$, so some value is hit $P^{2^k-k}$ times. Hua's lemma says that, on average, the sums of $2^{k-1}$ $k$-th powers are as well spread as they can be.

\section{Waring's problem}

Let $r_{s,k}(n)$ be the number of $(x_1,\dots,x_s)$ in positive integers with $x_1^k+\cdots+x_s^k=n$. Take $P=n^{1/k}$; then every solution has $x_i\le P$ and $r_{s,k}(n)=\int_0^1f(\alpha)^se(-n\alpha)\dd\alpha$. The major arc approximation involves the complete sums
\[
S(q,a)=\sum_{r=1}^qe\Bigl(\frac{ar^k}q\Bigr),\qquad v(\beta)=\int_0^Pe(\beta t^k)\dd t.
\]

\begin{lemma}\label{lem:waringmajor}
If $1\le q\le P$ and $\alpha=\frac aq+\beta$, then $f(\alpha)=q^{-1}S(q,a)v(\beta)+O_k(q(1+P^k|\beta|))$.
\end{lemma}
\begin{proof}
Since $x^k\equiv r^k\pmod q$ when $x\equiv r$,
\[
f(\alpha)=\sum_{r=1}^qe\Bigl(\frac{ar^k}q\Bigr)\sum_{\substack{x\le P\\x\equiv r}}g(x),\qquad g(t)=e(\beta t^k).
\]
 For fixed $r$, write $x=r+qm$, $0\le m\le M$. Each term $g(r+qm)$ differs from $\frac1q\int_{r+qm}^{r+q(m+1)}g$ by at most $\int_{r+qm}^{r+q(m+1)}|g'|$, and $\frac1q\int_{r}^{r+q(M+1)}g$ differs from $\frac1qv(\beta)$ by at most $2$, the two pieces $[0,r]$ and $[P,r+q(M+1)]$ having length at most $q$. With $|g'(t)|=2\pi k|\beta|t^{k-1}$ and $q\le P$, $\int_0^{2P}|g'|=2\pi|\beta|(2P)^k$. So the inner sum is $\frac{v(\beta)}q+O(1+P^k|\beta|)$; multiply by the $q$ roots of unity and add.
\end{proof}

\begin{lemma}\label{lem:singint}
(a) $|v(\beta)|\le\min(P,2|\beta|^{-1/k})$.
(b) If $s\ge2k$, then $\displaystyle\int_\R v(\beta)^se(-n\beta)\dd\beta=\frac{\Gamma(1+\frac1k)^s}{\Gamma(\frac sk)}\,n^{\frac sk-1}$.
\end{lemma}
\begin{proof}
Substituting $u=t^k$, $v(\beta)=\int_0^ng(u)e(\beta u)\dd u$ with $g(u)=\frac1ku^{1/k-1}$ on $\left]0,n\right]$ and $g=0$ elsewhere.
(a) $|v|\le P$ is clear. For $|\beta|\ge\frac1n$ let $U=\frac1{|\beta|}$; then $\int_0^Ug=U^{1/k}$. On $[U,n]$, $g$ is positive and decreasing, and $E(u)=\int_U^ue(\beta w)\dd w$ has $|E|\le\frac1{\pi|\beta|}$; integrating by parts, $|\int_U^nge(\beta u)\dd u|\le\frac{2g(U)}{\pi|\beta|}=\frac{2U^{1/k}}{k\pi}$. So $|v(\beta)|\le(1+\frac2{k\pi})U^{1/k}\le2|\beta|^{-1/k}$; for $|\beta|<\frac1n$, $P<|\beta|^{-1/k}$.
(b) By (a), $|v|^s$ is integrable when $s>k$. Let $g^{*j}$ be the $j$-fold convolution of $g$. The Beta integral of Chapter~7 gives, by induction on $j$, $\int_0^m\frac{u^{a-1}}{\Gamma(a)}\frac{(m-u)^{b-1}}{\Gamma(b)}\dd u=\frac{m^{a+b-1}}{\Gamma(a+b)}$, hence
\[
g^{*j}(m)=\frac{\Gamma(\frac1k)^j}{k^j\Gamma(\frac jk)}m^{\frac jk-1}=\frac{\Gamma(1+\frac1k)^j}{\Gamma(\frac jk)}m^{\frac jk-1}\qquad(0<m\le n),
\]
because for $m\le n$ the cut-off at $n$ plays no role; for $m>n$, $g^{*j}(m)$ is at most the same expression. In particular $g^{*k}$ is bounded, so $g^{*2k}=g^{*k}*g^{*k}$ is continuous (a convolution of two functions in $L^2$), and so is $g^{*s}=g^{*2k}*g^{*(s-2k)}$ when $s>2k$, a convolution of a bounded uniformly continuous function with an integrable one. Its Fourier transform is $\hat g(\xi)^s=v(-\xi)^s$, which is integrable. By the Fourier inversion of Chapter~10, at the point $n$, $g^{*s}(n)=\int v(-\xi)^se(n\xi)\dd\xi=\int v(\beta)^se(-n\beta)\dd\beta$.
\end{proof}

The condition $s\ge2k$ in (b) only serves to make the continuity easy; it holds in the theorem below, since $2^k+1\ge2k$.

\begin{theorem}[Hardy and Littlewood, 1920; Hua, 1938]\label{thm:waring}
Let $k\ge2$ and $s\ge2^k+1$. The series
\[
\mathfrak S_{s,k}(n)=\sum_{q\ge1}\sum_{\substack{a\bmod q\\(a,q)=1}}\Bigl(\frac{S(q,a)}q\Bigr)^se\Bigl(-\frac{an}q\Bigr)
\]
converges absolutely, uniformly in $n$, and there is $\delta=\delta(k)>0$ such that
\[
r_{s,k}(n)=\frac{\Gamma(1+\frac1k)^s}{\Gamma(\frac sk)}\,\mathfrak S_{s,k}(n)\,n^{\frac sk-1}+O_{s,k}\bigl(n^{\frac sk-1-\delta}\bigr).
\]
\end{theorem}
\begin{proof}
Fix $\varepsilon>0$ small in terms of $s$ and $k$, and let $\nu=\frac1{10}$.

\emph{Complete sums.} For $(a,q)=1$, $S(q,a)$ is the Weyl sum of length $q$ at $\alpha=\frac aq$, so Theorem~\ref{thm:weylk} with $P=q$ gives $|S(q,a)|\ll q^{1+\varepsilon}(\frac2q+q^{1-k})^{2^{1-k}}\ll q^{1-2^{1-k}+\varepsilon}$. As $s2^{1-k}\ge2+2^{1-k}$, the inner sum in $\mathfrak S_{s,k}$ is $O(q\cdot q^{-s2^{1-k}+s\varepsilon})=O(q^{-1-2^{1-k}+s\varepsilon})$, so the series converges absolutely, and its tail beyond $Q$ is $O(Q^{-2^{1-k}+s\varepsilon})$.

\emph{Arcs.} For $q\le P^\nu$ and $(a,q)=1$ let $\mathfrak M(q,a)=\{\alpha:|\alpha-\frac aq|\le P^{\nu-k}\}$; they are disjoint for $P$ large, as $P^{-2\nu}>2P^{\nu-k}$, and their union $\mathfrak M$ has measure at most $2P^{3\nu-k}$. On the complement $\mathfrak m$, Lemma~\ref{lem:dirichlet} with $Q=P^{k-\nu}$ gives $q\le Q$ with $|\alpha-\frac aq|\le\frac1{qQ}\le\frac1{q^2}$, and $q>P^\nu$ (otherwise $\alpha\in\mathfrak M$). Theorem~\ref{thm:weylk} then gives $|f(\alpha)|\ll P^{1+\varepsilon}(3P^{-\nu})^{2^{1-k}}$.

\emph{Minor arcs.} By Hua's lemma with $j=k$,
\[
\int_{\mathfrak m}|f|^s\le\sup_{\mathfrak m}|f|^{s-2^k}\int_0^1|f|^{2^k}\ll P^{(s-2^k)(1-\nu2^{1-k}+\varepsilon)}P^{2^k-k+\varepsilon}\le P^{s-k-\nu2^{1-k}+s\varepsilon}.
\]

\emph{Major arcs.} On $\mathfrak M(q,a)$, with $V=q^{-1}S(q,a)v(\beta)$, Lemma~\ref{lem:waringmajor} gives $|f-V|\ll P^\nu(1+P^\nu)\ll P^{2\nu}$, and $|f|,|V|\le P$, so $|f^s-V^s|\le sP^{s-1}|f-V|\ll P^{s-1+2\nu}$. Over $\mathfrak M$ this costs $O(P^{s-k-1+5\nu})$. The main part is
\[
\sum_{q\le P^\nu}\sum_{(a,q)=1}\Bigl(\frac{S(q,a)}q\Bigr)^se\Bigl(-\frac{an}q\Bigr)\int_{|\beta|\le P^{\nu-k}}v(\beta)^se(-n\beta)\dd\beta.
\]
By Lemma~\ref{lem:singint}(a), the integral over $|\beta|>P^{\nu-k}$ is at most $2\int_{P^{\nu-k}}^\infty(2\beta^{-1/k})^s\dd\beta\ll P^{(\nu-k)(1-s/k)}=P^{s-k-\nu(s/k-1)}$, so the integral is $J(n)+O(P^{s-k-\nu(s-k)/k})$ with $J(n)=\frac{\Gamma(1+1/k)^s}{\Gamma(s/k)}P^{s-k}$ by Lemma~\ref{lem:singint}(b). The sum over $q\le P^\nu$ is bounded, and differs from $\mathfrak S_{s,k}(n)$ by $O(P^{-\nu(2^{1-k}-s\varepsilon)})$.

\emph{Conclusion.} Every error is $O(P^{s-k-\delta'})$ for some $\delta'>0$ once $\varepsilon$ is small: the exponents saved are $\nu2^{1-k}-s\varepsilon$, $1-5\nu$, $\nu\frac{s-k}k$ and $\nu(2^{1-k}-s\varepsilon)$. Since $P^{s-k}=n^{s/k-1}$, the theorem holds with $\delta=\frac{\delta'}k$.
\end{proof}

The singular series is a product of local densities, $\mathfrak S_{s,k}(n)=\prod_p\chi_p(n)$, where $\chi_p(n)=\lim_{j\to\infty}p^{j(1-s)}\#\{\mathbf x\bmod p^j:x_1^k+\cdots+x_s^k\equiv n\}$ measures how often $n$ is a sum of $s$ $k$-th powers $p$-adically. It can vanish: fourth powers are $0$ or $1$ modulo $16$, so for $s\le14$ no $n\equiv15\pmod{16}$ is a sum of $s$ fourth powers even modulo $16$, and $\chi_2(n)=0$. Hensel's lemma shows that it is bounded below as soon as there are enough variables.

\begin{theorem}[Hardy and Littlewood; quoted]\label{thm:singpos}
If $s\ge4k$, then $\mathfrak S_{s,k}(n)\ge c(s,k)>0$ for all $n\ge1$.
\end{theorem}

\begin{corollary}[Hua, 1938]\label{cor:gk}
For $k\ge4$, every sufficiently large integer is a sum of $2^k+1$ positive $k$-th powers: $G(k)\le2^k+1$.
\end{corollary}
\begin{proof}
For $k\ge4$, $s=2^k+1\ge4k$, so Theorems~\ref{thm:waring} and~\ref{thm:singpos} give $r_{s,k}(n)\ge c\,n^{s/k-1}-O(n^{s/k-1-\delta})>0$ for $n$ large.
\end{proof}

\begin{figure}[ht]
\centering
\fig{ch07-fivesquares}
\caption{Sums of five squares ($k=2$, $s=5=2^k+1$), for $n=2000,2010,\dots,4000$. Horizontally, the singular series $\mathfrak S_{5,2}(n)$, summed over $q\le200$; vertically, $r_{5,2}(n)$, computed exactly, divided by $\frac{\Gamma(3/2)^5}{\Gamma(5/2)}n^{3/2}$. The correlation is $0.993$. The points lie along a line slightly below the diagonal (the average ratio is $0.93$), because of a lower-order term of relative size about $n^{-1/2}$, coming from the representations that would use $0$.}
\label{fig:fivesquares}
\end{figure}

\begin{beauty}
A count of solutions is written as an integral over the circle. Near the fractions with small denominators the integrand factorizes into a \emph{local} part, a complete exponential sum modulo $q$, and an \emph{archimedean} part, an integral over the reals. Summing over all $q$ assembles the local parts into a product over all primes. The main term is then a product of densities, one for each completion of $\Q$: the singular integral for $\R$, the factor $\chi_p$ for each $\Q_p$. Everything else, the minor arcs, is noise, and the whole difficulty is to prove that it is noise. Volume~14 runs the same argument for three primes.
\end{beauty}

\section{Exercises}

\begin{exercise}[$\star$]\label{ex:gausssq}
For $k=2$ and $(a,q)=1$, show that $|S(q,a)|^2\le2q$. Deduce directly, without Weyl's inequality, that $\mathfrak S_{s,2}(n)$ converges absolutely for $s\ge5$.
\end{exercise}
\begin{exercise}[$\star\star$]\label{ex:density}
Use Hua's lemma and Cauchy--Schwarz to show that at least $X^{1-\varepsilon}$ of the integers $n\le X$ are sums of $2^{k-1}$ positive $k$-th powers, for $X\ge X_0(k,\varepsilon)$. What does it say for four cubes?
\end{exercise}
\section{Hints}

\begin{hint}{ex:gausssq}
Expand $|S(q,a)|^2$, put $x=y+h$, and sum over $y$ first: $x^2-y^2=h^2+2hy$.
\end{hint}
\begin{hint}{ex:density}
$\sum_nr_{s,k}(n)=\lfloor P\rfloor^s$ and $\sum_nr_{s,k}(n)^2=\int_0^1|f|^{2s}$; take $2s=2^k$.
\end{hint}
\section{Solutions}

\begin{solution}{ex:gausssq}
\idea Squaring a Gauss sum linearizes the phase in one of the two variables.
\writeup $|S(q,a)|^2=\sum_{x,y\bmod q}e(\frac{a(x^2-y^2)}q)$. With $x=y+h$, $x^2-y^2=h^2+2hy$, so
\[
|S(q,a)|^2=\sum_{h\bmod q}e\Bigl(\frac{ah^2}q\Bigr)\sum_{y\bmod q}e\Bigl(\frac{2ahy}q\Bigr).
\]
The inner sum is $q$ if $q\mid2ah$, that is $q\mid2h$ since $(a,q)=1$, and $0$ otherwise. At most two classes $h$ modulo $q$ have $q\mid2h$, so $|S(q,a)|^2\le2q$. Then the $q$-th term of $\mathfrak S_{s,2}$ is at most $q(\frac{2}{q})^{s/2}=2^{s/2}q^{1-s/2}$, and $\sum_qq^{1-s/2}$ converges for $s\ge5$.
\end{solution}
\begin{solution}{ex:density}
\idea If the values were concentrated on few integers, the number of coincidences, which is the $2s$-th moment, would be large.
\writeup Let $s=2^{k-1}$, $P=X^{1/k}$, and let $Z$ be the number of $n\le X$ with $r_{s,k}(n)>0$; all values $\sum x_i^k$ with $x_i\le P/s^{1/k}$ are at most $X$, so with $P'=P/s^{1/k}$ and $r'$ counting those, Cauchy--Schwarz gives
\[
\lfloor P'\rfloor^{2s}=\Bigl(\sum_nr'(n)\Bigr)^2\le Z\sum_nr'(n)^2=Z\int_0^1|f_{P'}|^{2^k}\ll ZP'^{2^k-k+\varepsilon}
\]
by Hua's lemma with $j=k$ ($2s=2^k$). So $Z\gg P'^{k-\varepsilon}\gg X^{1-\varepsilon/k}$. For cubes, $s=4$: almost every integer, in the sense that $X^{1-\varepsilon}$ of those up to $X$, is a sum of four positive cubes. Much more is true: Davenport proved in 1939, with a refined circle method, that almost all integers are sums of four positive cubes.
\end{solution}
\section*{Sources and further reading}

\begin{itemize}
\item \emph{History.} E.~Waring, \emph{Meditationes Algebraicae}, Cambridge, 1770. D.~Hilbert, ``Beweis für die Darstellbarkeit der ganzen Zahlen durch eine feste Anzahl $n^{\text{ter}}$ Potenzen (Waringsches Problem)'', \emph{Math.\ Ann.}\ 67 (1909), 281--300. On Hua: J.~J.~O'Connor and E.~F.~Robertson, ``Hua Loo-Keng'', MacTutor History of Mathematics archive.
\item \emph{The circle method.} G.~H.~Hardy and J.~E.~Littlewood, ``Some problems of `Partitio Numerorum' I: A new solution of Waring's problem'', \emph{Nachr.\ Ges.\ Wiss.\ Göttingen} (1920), 33--54. L.-K.~Hua, ``On Waring's problem'', \emph{Quart.\ J.\ Math.\ Oxford} 9 (1938), 199--202. R.~C.~Vaughan, \emph{The Hardy--Littlewood Method}, 2nd ed., Cambridge, 1997, Chapters~2--4 (Theorem~\ref{thm:singpos} is Theorem~4.6 there). H.~Davenport, \emph{Analytic Methods for Diophantine Equations and Diophantine Inequalities}, 2nd ed., Cambridge, 2005.
\end{itemize}
